% T33 — «Письмо ученику»: стилизация под рукописное письмо.
% Засечный italic, без рамок, задачи как «Пункт 1.», «Пункт 2.». Math, класс 11.
\documentclass[a4paper,12pt]{article}

\usepackage{../../../../Lessons/_templates/preamble_worksheet}
\usepackage{../_shared/worksheet_check}

\geometry{a4paper, margin=22mm}

\usepackage{eso-pic}

% --- Палитра: чернильный синий + кремовая бумага ---
\definecolor{wcprimary}{HTML}{1B2D5C}
\definecolor{wcprimaryLight}{HTML}{C1C8DC}
\definecolor{wcprimaryBG}{HTML}{FBF7EE}
\definecolor{wcink}{HTML}{2A3D6D}

% Шрифт оставим Times (засечки)

% --- Кремовый фон ---
\AddToShipoutPictureBG{%
  \AtPageLowerLeft{%
    \color{wcprimaryBG}\rule{\paperwidth}{\paperheight}}}

% --- Минимальный заголовок: только подпись и дата ---
\renewcommand{\WorksheetTitle}[1]{%
  \par\noindent
  {\Large\itshape\color{wcink} #1}\par
  \vspace{1mm}
  {\color{wcink}\rule{30mm}{0.4pt}}\par
  \vspace{4mm}}

% Пункт письма
\newcommand{\LetterPara}[2]{% #1 номер #2 содержание
  \par\vspace{2mm}
  \noindent
  \textbf{\itshape\color{wcink}Пункт #1.}\ #2
  \par}

\SetAnswerFieldSize{36mm}{8mm}

% Переопределим FirstPageHeader, чтобы без рамки
\renewcommand{\FirstPageHeader}{%
  \par\noindent
  {\itshape\color{wcink}\rule{0pt}{4mm}Москва, \today}\hfill
  {\itshape\color{wcink}кому: \underline{\hspace{50mm}}}\par
  \vspace{4mm}}

% QR без рамок будет сам по себе

\begin{document}

\WorksheetTitle{Дорогой ученик!}
\FirstPageHeader

\noindent
\itshape
Сегодня я хочу поделиться с тобой пятью задачами, которые помогут вспомнить
самое важное из курса алгебры и начала анализа. Не торопись~--- читай,
размышляй, и пусть каждый пункт станет небольшим открытием.\par
\upshape
\vspace{2mm}

\LetterPara{1}{%
  \itshape
  Вспомни производную. Найди $f'(x)$, если $f(x) = x^3 - 6x^2 + 9x + 2$,
  и определи, при каком целом $x$ функция достигает локального максимума.\par
  \upshape\vspace{1mm}
  \hfill\textbf{\color{wcink}Ответ:}~\AnswerField{1}%
}

\LetterPara{2}{%
  \itshape
  Логарифмы~--- это просто, если помнить их свойства. Реши уравнение
  $\log_2(x-1) + \log_2(x+3) = 5$. Подумай, какое из решений подходит по~ОДЗ.\par
  \upshape\vspace{1mm}
  \hfill\textbf{\color{wcink}Ответ:}~\AnswerField{2}%
}

\LetterPara{3}{%
  \itshape
  Тригонометрия. Найди наибольшее значение функции $y = 4\sin x - 3\cos x$.
  Помни про вспомогательный угол.\par
  \upshape\vspace{1mm}
  \hfill\textbf{\color{wcink}Ответ:}~\AnswerField{3}%
}

\LetterPara{4}{%
  \itshape
  Геометрическая задача, которая ждёт твоего внимания: в правильной
  четырёхугольной пирамиде сторона основания равна $6$, а высота равна $4$.
  Найди объём пирамиды.\par
  \upshape\vspace{1mm}
  \hfill\textbf{\color{wcink}Ответ:}~\AnswerField{4}%
}

\LetterPara{5}{%
  \itshape
  И напоследок~--- задача с параметром. При каком значении $a$ уравнение
  $x^2 - 2ax + a + 6 = 0$ имеет ровно один корень? Запиши их сумму.\par
  \upshape\vspace{1mm}
  \hfill\textbf{\color{wcink}Ответ:}~\AnswerField{5}%
}

\vspace{4mm}
\noindent
\itshape
Удачи тебе! Помни: математика не наказание, а способ навести порядок в мыслях.\par
\vspace{2mm}
\hfill {\itshape\color{wcink}Твой учитель.}
\upshape
\par

\WorksheetQR{T33-letter/L1/v1}

\end{document}
